\chapter{Discussion}
\label{ch:discussion}

\section{Answering the research question}

The research question was: \emph{how do recurrent networks and Transformers compare at classifying
short news texts?} The evidence of Chapter~\ref{ch:results} supports a four-part answer.

\subsection{They are indistinguishable in accuracy}

The measured difference is 0.09 percentage points. McNemar's exact test on the paired predictions
gives $p = 0.754$, the models agree on 94.78\% of test documents, and the discordant set splits 187
to 180. There is no basis in this evidence for preferring either architecture on accuracy grounds.

This is a negative result, and it is reported as the principal finding rather than buried. Negative
results of this kind are underreported in the literature, with the consequence that accuracy
differences well inside the noise floor are routinely presented as architectural improvements. The
discipline of applying a paired significance test is what converts an apparent 0.09-point win into
the correct conclusion of equivalence.

\subsection{The explanation is document length}

The equivalence is not a coincidence, and it is not evidence that self-attention is overrated in
general. It follows from a specific property of the data.

Self-attention's structural advantage, established in Table~\ref{tab:complexity}, is a maximum path
length of $O(1)$ between any two positions, against $O(n)$ for recurrence. This is an advantage in
modelling \emph{long-range} dependency.

The median AG News document is 39 tokens long. In a 39-token document there are essentially no
long-range dependencies to model. The furthest two tokens can be is a few dozen positions, which a
bidirectional LSTM traverses without difficulty --- this is precisely the regime the gating
mechanism of Section~\ref{sec:vanishing} was designed to handle, and it is two orders of magnitude
short of where the gating mechanism begins to fail.

Furthermore, the signal in topic classification is overwhelmingly lexical. The presence of
\emph{coach}, \emph{Olympic} and \emph{medal} identifies a sports story regardless of their
arrangement. This observation is corroborated independently by the literature: FastText, which
averages word vectors and models no order whatsoever, reaches approximately 92.5\% on this
benchmark (Table~\ref{tab:prior}) --- at or above both models in this study. If a bag-of-embeddings
model is competitive, then the marginal value of \emph{any} sequence modelling on this task is
small, and the marginal value of one sequence model over another is smaller still.

\keypoint{The general form of the argument.} An architecture's advantage can only be realised on a
task that exercises the property the architecture is good at. Transformers are superior at modelling
long-range dependency; AG News contains almost none; therefore the two architectures perform
equivalently. This is not a limitation of the Transformer but a statement about the task.

\subsection{They fail in the same way}

The equivalence runs deeper than aggregate accuracy. The two models share:

\begin{itemize}
\item an identical per-class difficulty ordering (Sports, World, Sci/Tech, Business),
\item near-identical per-class $F_1$, agreeing to four decimal places on Business,
\item the same dominant confusion, Business predicted as Sci/Tech, at 9.2\% and 9.6\%,
\item agreement on 94.78\% of individual test predictions.
\end{itemize}

Two models with materially different inductive biases would be expected to differ somewhere in this
profile even at equal overall accuracy. That they do not indicates both are limited by the same
external factor. The error analysis of Chapter~\ref{ch:results} identifies that factor: documents
whose gold label is one of several admissible choices. A story about a semiconductor firm's
quarterly earnings is legitimately Business and legitimately Sci/Tech; the corpus assigns one
label, and no model can recover the other.

This is an irreducible component of the error, a property of the annotation rather than of either
model, and it caps what any architecture can achieve on this benchmark.

\subsection{Cost is where they differ, and the finding inverts expectation}

At matched accuracy and matched capacity, the recurrent model trained 1.67 times faster per epoch
and performed inference 2.43 times faster.

Section~\ref{sec:cost} analyses why. The summary is that the Transformer's advantage is
\emph{asymptotic} --- it concerns parallelisability and the growth of sequential depth with
sequence length --- and that neither condition for realising it holds in this experiment. There is
negligible hardware parallelism to exploit on six CPU threads, and at $n = 64$, $d = 128$ the
$O(n d^2)$ and $O(n^2 d)$ arithmetic costs are of the same order, so self-attention does not even
win on operation count.

The honest framing is that this is a result about a \emph{regime}, not about the architectures in
the abstract. Two predictions follow, and both are testable: the ordering should reverse on
parallel hardware, and the recurrent model's relative cost should grow linearly with sequence
length.

\section{Comparison with published results}

Both models fall within the 90--92\% band that Table~\ref{tab:prior} identifies for non-pretrained
approaches on AG News. The study therefore reproduces the expected performance level, which
provides weak corroboration that the pipeline as a whole is sound.

The roughly three-point gap to fine-tuned BERT is instructive and is worth stating plainly, since
it is frequently misattributed. That gap is \textbf{not} evidence that self-attention is a superior
architecture. It is evidence that representations learned from billions of words of unlabelled text
are superior to representations learned from 54{,}000 labelled headlines. The present study
isolates architecture by holding pretraining constant at zero for both models, and under that
control the architectural difference vanishes entirely.

This distinction matters practically. An engineer choosing between a BiLSTM and a Transformer
encoder for a short-text classification task, both trained from scratch, should not expect the
accuracy gain reported for BERT to transfer to that choice.

\section{Implications}

\keypoint{For practice.} On short-text topic classification at this scale, selecting between these
two architectures on expected accuracy is not supported by evidence. The decision should be made on
engineering grounds --- latency, memory, deployment constraints, tooling --- and those costs should
be \emph{measured on the target hardware and sequence length} rather than assumed from asymptotic
arguments. The cost finding here is a concrete demonstration that the assumed ordering can be
wrong.

\keypoint{For evaluation methodology.} Reporting a sub-point accuracy difference without a paired
significance test is not sound practice. The test applied here cost a few lines of code and
reversed the qualitative conclusion, from ``the BiLSTM is slightly better'' to ``the models are
indistinguishable''.

\keypoint{For from-scratch implementation.} The verification suite of
Chapter~\ref{ch:verification} is a transferable methodology. The key insight is that numerical code
of this kind fails silently, so verification must consist of equivalence against independent
implementations and assertion of mathematical invariants, not merely the observation that training
loss decreases.

\section{Threats to validity}
\label{sec:threats}

\subsection{Internal validity}

\keypoint{Single seed.} This is the most serious limitation. Each architecture was trained once,
with one seed. The observed 0.09-point difference is well within typical run-to-run variance for
models of this size, which is precisely why Section~\ref{sec:significance} reports a significance
test rather than a winner. The conclusion supported is ``indistinguishable''; a claim that one
architecture is genuinely superior --- in either direction --- would require multiple seeds and a
confidence interval. The study does not make such a claim.

\keypoint{Learning rates differ.} Defended in Chapter~\ref{ch:method}: the controlled quantity is
the compute budget, and both rates were selected on validation data. Residual risk remains that a
more exhaustive search would favour one architecture.

\keypoint{Both models overfit quickly.} The recurrent model peaked after one epoch. Stronger
regularisation might alter the comparison, although both models were selected at their respective
optima by the early-stopping criterion.

\subsection{External validity}

\keypoint{One corpus, one length regime.} All conclusions concern English short news text with a
median length of 39 tokens. The central finding is explicitly \emph{conditional} on that length,
and the report predicts it would not hold for longer documents.

\keypoint{No pretraining.} The findings concern architectures trained from scratch and do not
transfer to the pretrained setting.

\keypoint{CPU-only timing.} Every figure in Section~\ref{sec:cost} is specific to the hardware of
Table~\ref{tab:environment} and would be expected to change on a GPU.

\keypoint{Reduced training set.} 60{,}000 of 120{,}000 available records were used, per the brief.
The more parameter-hungry Transformer might benefit disproportionately from the full corpus.

\subsection{Construct validity}

\keypoint{Truncation.} \code{MAX\us{}LEN} of 64 truncates 3.06\% of training documents. This affects
both models identically and is unlikely to bias the comparison.

\keypoint{Simple tokenisation.} A subword tokeniser would likely raise both models' absolute
accuracy. It was omitted to keep the input identical across architectures, at a known cost in
absolute performance.

\keypoint{Label noise.} As the error analysis shows, a meaningful proportion of test documents
admit more than one defensible label. This places a ceiling on measurable accuracy and compresses
the range within which architectures can be distinguished --- which, it should be acknowledged,
makes an equivalence finding easier to obtain. A task with cleaner labels and longer documents
would be a more discriminating test.
